De maneira aditiva ao resultado deste trabalho, 
a aplicação de métricas de redes de comunicação auxiliará na comparação entre ambos os métodos. 
Tais métricas foram selecionadas com base nos trabalhos de \citeonline{bravo2025ants} e \citeonline{}, 
que tem por foco o levantamento deste tipo de métrica em sistemas \gls{iot}.

\paragraph{Razão de Pacotes Entregues}
A robustez da transmissão via \gls{lorawan} será quantificada através da \gls{pdr}.  
O PDR é definido pela razão entre o número de pacotes recebidos 
pelo \textit{gateway} e o número de pacotes enviados pelo nó sensor:

\begin{equation}
\label{eq:pdr}
PDR = \left( \frac{\text{Pacotes Recebidos}}{\text{Pacotes Enviados}} \right) \times 100
\end{equation}


\paragraph{Latência e entrega}
\label{subsubsec:metricas_latencia}
A latência ponta a ponta de uma leitura é o intervalo entre
o disparo no nó e o instante em que o dado fica disponível no servidor: a imagem remontada,
no fluxo em nuvem, ou a leitura decodificada, no fluxo em borda. 

\begin{equation}
    \label{eq:entrega_leitura}
    E_{\mathrm{leitura}} = \frac{\text{leituras disponíveis no servidor}}{\text{leituras disparadas}}
\end{equation}

\paragraph{Avaliação}

As duas métricas diferem no fluxo em nuvem: como nenhum dos firmwares retransmite blocos no
modo autônomo, uma imagem com um único bloco perdido não é remontada, e a leitura se perde
mesmo com um \gls{pdr} alto. No fluxo em borda, em que cada leitura ocupa um pacote, as duas
coincidem.

